\documentclass[11pt,a4paper]{article}
\usepackage{templates/modeling-notes}
\usepackage{needspace}
\hypersetup{pdftitle={Chapter 3 -- The simplex method}}
\title{\textbf{Mathematical Modeling}\\[2mm]
\large Chapter 3: The Simplex Method}
\author{Personal study notes}
\date{}
\begin{document}
\maketitle
{\small\tableofcontents}

\section{The problem and its possible outcomes}\label{sec:problem}
The primal simplex method solves a linear program by maintaining a feasible
basis and exchanging one basic variable at a time. For a maximization problem,
use the standard form
\begin{equation}\label{eq:standard}
  \max c^Tx,\qquad Ax=b,\qquad x\geq0,
  \quad A\in\R^{m\times n},\ b\in\R^m,\ c,x\in\R^n.
\end{equation}
Unless explicitly discussing redundant rows, assume $\operatorname{rank}A=m$
and $m\leq n$. If $m=n$, the nonsingular system has a single candidate
$x=A^{-1}b$: check its nonnegativity directly. The interesting basis exchanges
occur when $m<n$. Minimization can be handled by maximizing the negative
objective and reversing the sign of the final value.

There are three main outcomes. An infeasible problem has no point satisfying
all constraints. A feasible problem with objective unbounded above has no finite
maximum. Otherwise the maximum is finite and attained; in standard form at
least one optimal solution is a basic feasible solution. The optimal solution
may be unique or there may be infinitely many: the segment between two distinct
optimal points is also optimal.
\begin{figure}[H]
  \centering
  \begin{tikzpicture}[every node/.style={font=\small},
      box/.style={draw=modelblue,rounded corners,align=center,inner sep=6pt}]
    \node[box] (lp) at (0,0) {Linear program};
    \node[box] (empty) at (-4,-1.3) {Infeasible\\$\Omega=\varnothing$};
    \node[box] (feasible) at (2,-1.3) {Feasible\\$\Omega\ne\varnothing$};
    \node[box] (finite) at (0,-2.7) {Finite maximum\\attained};
    \node[box] (unbounded) at (4.2,-2.7) {Unbounded objective\\$\sup c^Tx=+\infty$};
    \node[box] (one) at (-1.7,-4) {Unique optimizer};
    \node[box] (many) at (1.7,-4) {Infinitely many\\optimizers};
    \foreach \a/\b in {lp/empty,lp/feasible,feasible/finite,feasible/unbounded,finite/one,finite/many}
    \draw[model flow] (\a)--(\b);
  \end{tikzpicture}
  \caption{Classify feasibility first, then objective behavior. An unbounded
    feasible set can still have a finite maximum.}
  \label{fig:outcomes}
\end{figure}
For example, $\max(-x)$ over $x\geq0$ has an unbounded feasible set but the
unique optimizer $x=0$. Unboundedness of the set and of the objective are
separate questions.

\Needspace{14cm}
\section{Bases, basic solutions, and degeneracy}\label{sec:bases}
Let $B=(j_1,\ldots,j_m)$ be an ordered list of indices whose columns form the
nonsingular matrix $A_B=[a_{j_1}\ \cdots\ a_{j_m}]$. Let $N$ contain the
remaining indices, and partition variables and objective coefficients in the
same order. Then
\[
  A_Bx_B+A_Nx_N=b,
  \qquad x_B=A_B^{-1}b-A_B^{-1}A_Nx_N.
\]
The basic solution associated with $B$ sets $x_N=0$ and solves
$A_Bx_B=b$. It is a \emph{basic feasible solution} (BFS, Polish BRD) exactly
when $x_B\geq0$. A feasible basic solution is \emph{degenerate} when at least
one basic coordinate is zero. A nonbasic zero does not establish degeneracy.

\begin{figure}[H]
  \centering
  \begin{tikzpicture}[every node/.style={font=\small}]
    \node at (0,1.4) {Choose $\color{modelblue}B=(2,4)$;
      the remaining indices are $\color{modelorange}N=(1,3)$.};
    \node at (0,0) {$
        \underbrace{\begin{bmatrix}
            \color{modelorange}1 & \color{modelblue}1 & \color{modelorange}2 & \color{modelblue}0\\
            \color{modelorange}1 & \color{modelblue}0 & \color{modelorange}1 & \color{modelblue}1
          \end{bmatrix}}_{A}
        \underbrace{\begin{bmatrix}
            \color{modelorange}x_1\\ \color{modelblue}x_2\\
            \color{modelorange}x_3\\ \color{modelblue}x_4
          \end{bmatrix}}_{x}
        =\underbrace{\begin{bmatrix}4\\3\end{bmatrix}}_{b}
      $};
    \draw[model flow] (0,-1.3)--(0,-2.1);
    \node[right,align=left] at (0.3,-1.7) {Group columns and variables\\in the same order.};
    \node at (0,-3.3) {$
      \underbrace{\left[\begin{array}{cc|cc}
            \color{modelblue}1 & \color{modelblue}0 & \color{modelorange}1 & \color{modelorange}2\\
            \color{modelblue}0 & \color{modelblue}1 & \color{modelorange}1 & \color{modelorange}1
          \end{array}\right]}_{[\,\color{modelblue}A_B\ \color{modelorange}A_N\,]}
      \underbrace{\left[\begin{array}{c}
            \color{modelblue}x_2\\ \color{modelblue}x_4\\\hline
            \color{modelorange}x_1\\ \color{modelorange}x_3
          \end{array}\right]}_{\binom{\color{modelblue}x_B}{\color{modelorange}x_N}}
      =\underbrace{\begin{bmatrix}4\\3\end{bmatrix}}_{b}
    $};
    \node[draw=modelblue,rounded corners,inner sep=6pt] at (0,-5.7)
    {$x_N=\binom{x_1}{x_3}=0
        \quad\Longrightarrow\quad x_B=\binom{x_2}{x_4}=\binom{4}{3}
        \quad\Longrightarrow\quad x=\begin{pmatrix}0\\4\\0\\3\end{pmatrix}$.};
  \end{tikzpicture}
  \caption{An illustrative partition with $m=2$, $n=4$. Blue marks the basic
    columns and variables; orange marks the nonbasic ones. Here $A_B=I_2$,
    so setting $x_N=0$ immediately gives a nondegenerate BFS.
    Grouping changes the display order, not the original variable numbering.}
  \label{fig:partition}
\end{figure}
A BFS is a vertex of $\{x:Ax=b,x\geq0\}$, and conversely every vertex admits a
feasible basis. Different bases can represent the same degenerate vertex.
There are at most $\binom{n}{m}$ bases, but some column subsets are singular,
some basic solutions are infeasible, and some solutions coincide. This is an
upper bound, not an exact count of vertices.

\Needspace{22\baselineskip}
\subsection{Classifying three candidate solutions}
Consider
\begin{equation}\label{eq:classification}
  \begin{aligned}
    2x_1+x_2+3x_3+x_6&=8,\\
    x_1+2x_3+x_5+x_6&=5,\\
    3x_1+x_3+x_4+2x_6&=4.
  \end{aligned}
\end{equation}
The columns of $A$ and the right-hand side are:
\begin{center}
  \begin{tabular}{c*{6}{r}@{\qquad}r}\toprule
    Equation & $a_1$ & $a_2$ & $a_3$ & $a_4$ & $a_5$ & $a_6$ & $b$\\\midrule
    1 & 2 & 1 & 3 & 0 & 0 & 1 & 8\\
    2 & 1 & 0 & 2 & 0 & 1 & 1 & 5\\
    3 & 3 & 0 & 1 & 1 & 0 & 2 & 4\\\bottomrule
  \end{tabular}
\end{center}
Here $A$ has full row rank because columns $(2,5,4)$ form $I_3$.
Solving the indicated three-column systems gives:
\begin{center}
  \begin{tabular}{llll}\toprule
    Candidate $x^T$ & Ordered basis $B$ & Basic values $x_B^T$ & Classification\\\midrule
    $(0,8,0,4,5,0)$ & $(2,5,4)$ & $(8,5,4)$ & Nondegenerate BFS\\
    $(0,\frac12,\frac52,\frac32,0,0)$ & $(2,3,4)$ & $(\frac12,\frac52,\frac32)$ & Nondegenerate BFS\\
    $(0,0,0,-12,-3,8)$ & $(4,5,6)$ & $(-12,-3,8)$ & Basic, infeasible\\\bottomrule
  \end{tabular}
\end{center}
For the second candidate, $2x_3=5$, then $x_2+3x_3=8$ and $x_3+x_4=4$.
For the third, $x_6=8$, $x_5=5-8=-3$, and $x_4=4-16=-12$.
Each listed basis is nonsingular. Checking $Ax=b$ alone cannot establish
feasibility; conversely, counting nonzero entries alone cannot establish that
an arbitrary solution is basic without checking column independence.

\section{Canonical equations and reduced costs}\label{sec:canonical}
For a chosen feasible basis, define
\begin{equation}\label{eq:transformed}
  \bar A=A_B^{-1}A,\qquad \bar b=A_B^{-1}b\geq0.
\end{equation}
The basic columns of $\bar A$ form the identity. The resulting canonical
equations are
\[
  x_B+\bar A_Nx_N=\bar b.
\]
This is canonical \emph{relative to the basis}: any nonsingular basis produces
identity columns after transformation, but only a feasible basis gives
$\bar b\geq0$. In numerical code solve systems with $A_B$ rather than computing
its inverse explicitly.

Substitute the canonical equations into the objective:
\begin{align}
  c^Tx&=c_B^Tx_B+c_N^Tx_N\notag\\
  &=c_B^T\bar b+(c_N^T-c_B^T\bar A_N)x_N\notag\\
  &=z_0-\sum_{j\in N}q_jx_j,\qquad
  z_0=c_B^T\bar b,\quad q_j=c_B^T\bar a_j-c_j.\label{eq:reduced}
\end{align}
The indicators $q_j=z_j-c_j$ use the \emph{maximization} sign convention.
Another common convention calls $\rho_j=c_j-z_j=-q_j$ the reduced cost.
For every basic column, $q_j=0$.
\begin{center}
  \begin{tabular}{lll}\toprule
    Indicator & Equivalent reduced cost & Effect of increasing $x_j$\\\midrule
    $q_j<0$ & $\rho_j>0$ & Objective increases if a positive step is feasible\\
    $q_j=0$ & $\rho_j=0$ & Objective stays constant along this direction\\
    $q_j>0$ & $\rho_j<0$ & Objective decreases\\\bottomrule
  \end{tabular}
\end{center}

\begin{proposition}[Optimality certificate]\label{prop:certificate}
  If the basis is feasible and $q_j\geq0$ for every nonbasic index, then its
  basic solution is globally optimal.
\end{proposition}
\begin{proof}
  Every feasible point obeys the same transformed equations. Since its
  nonbasic coordinates are nonnegative, Equation~\eqref{eq:reduced} gives
  $c^Tx\leq z_0$. The current BFS attains $z_0$.
\end{proof}
For a nondegenerate BFS the condition is also necessary: if $q_k<0$, all
basic coordinates are positive, so a sufficiently small positive step in
$x_k$ preserves nonnegativity and strictly improves the objective.
At a degenerate BFS that step can be blocked at zero. For example,
$\max x_2$ subject to $x_1+x_2=0$, $x\geq0$ has just one feasible point,
$(0,0)$. For $B=(1)$, $q_2=-1$, yet the point is optimal. Changing to
$B=(2)$ gives $q_1=1$ and a certificate. Optimality of a point and the sign
certificate of a particular basis must be distinguished.

If all nonbasic $q_j>0$, equality in the certificate forces $x_N=0$ and hence
$x_B=\bar b$: the optimizer is unique. This is sufficient, not necessary.
If some nonbasic $q_k=0$, a \emph{positive feasible step} in that direction
produces distinct optimal points. A zero indicator by itself is insufficient:
$\max0$ with $x_1+x_2=0$, $x\geq0$ has $q_2=0$ and still only one feasible point.

\section{One simplex step}\label{sec:pivot}
\subsection{Entering variable and feasible direction}
Choose a nonbasic $k$ with $q_k<0$. Dantzig's rule chooses a most negative
indicator; any negative indicator is eligible. Hold other nonbasic coordinates
at zero and increase $x_k=\theta\geq0$. With $u=\bar a_k$,
\begin{equation}\label{eq:direction}
  x_B(\theta)=\bar b-\theta u,\qquad
  d_B=-u,\quad d_k=1,\quad d_j=0\ (j\in N\setminus\{k\}).
\end{equation}
Then $Ad=0$ and $c^Td=-q_k>0$. The equations remain satisfied along
$x(\theta)=x^{(B)}+\theta d$; nonnegativity determines the allowed step.

\subsection{Ratio test and leaving variable}
If $u_i>0$, basic coordinate $x_{j_i}$ decreases and requires
$\theta\leq\bar b_i/u_i$. If $u_i\leq0$, it places no upper bound on the step.
Thus
\begin{equation}\label{eq:ratio}
  \theta^*=\min_{i:u_i>0}\frac{\bar b_i}{u_i}.
\end{equation}
Choose a minimizing row $r$ and remove its basic variable $x_{j_r}$.
The pivot element is $p=\bar a_{rk}>0$. Replace $j_r$ by $k$ in the ordered
basis; $x_k=\theta^*$ and the leaving variable becomes nonbasic and zero.
The objective gain is $-q_k\theta^*$. A tie may leave another zero basic
coordinate, making the new BFS degenerate.

\begin{figure}[H]
  \centering
  \begin{tikzpicture}[x=0.85cm,y=0.42cm,every node/.style={font=\small}]
    \draw[model axis] (0,0)--(5.5,0) node[right] {$\theta$};
    \draw[model axis] (0,0)--(0,8.6) node[above] {basic value};
    \draw[modelblue,thick] (0,6)--(4,0) node[below] {$6/(3/2)=4$};
    \draw[modelorange,thick] (0,8)--(4,4);
    \draw[black!60,thick] (0,3)--(4,5);
    \draw[dashed] (4,0)--(4,7.8);
    \node[left] at (0,6) {$6$};
    \node[left] at (0,8) {$8$};
    \node[left] at (0,3) {$3$};
    \node[right] at (4,4) {$8-\theta$};
    \node[right] at (4,5) {$3+\theta/2$};
    \node[model point] at (4,0) {};
    \node[align=left] at (10,6.5) {Only decreasing coordinates\\limit the step.\\First zero determines the exit.};
  \end{tikzpicture}
  \caption{Illustrative ratio test: the bounds are $4$ and $8$; the increasing
    coordinate gives no bound. The maximum feasible step is $\theta^*=4$.}
  \label{fig:ratio}
\end{figure}

\subsection{When there is no leaving variable}
If $u\leq0$, the minimum in Equation~\eqref{eq:ratio} has no candidates.
Then $d\geq0$, so $x^{(B)}+\theta d$ is feasible for all $\theta\geq0$, and
$c^Tx(\theta)=z_0-\theta q_k\to+\infty$. This is a certificate of an
unbounded objective. It requires \emph{both} a negative indicator and a
column with no positive entry. The same column with $q_k=0$ instead gives
constant objective along a feasible ray.

\subsection{Tableau update by row operations}\label{sec:row-operations}
Write a tableau as the constraint rows $[\bar A\mid\bar b]$, followed by
$[q^T\mid z_0]$. The last row expresses
$z+\sum_jq_jx_j=z_0$. For a pivot in row $r$, column $k$, use old coefficients
on the right-hand sides of all update formulas:
\begin{align*}
  \bar a'_{rj}&=\bar a_{rj}/p,& \bar b'_r&=\bar b_r/p,\\
  \bar a'_{ij}&=\bar a_{ij}-\bar a_{ik}\bar a_{rj}/p,&
  \bar b'_i&=\bar b_i-\bar a_{ik}\bar b_r/p\quad(i\ne r),\\
  q'_j&=q_j-q_k\bar a_{rj}/p,& z'_0&=z_0-q_k\bar b_r/p.
\end{align*}
Equivalently: divide the pivot row by $p$, then subtract suitable multiples of
that normalized row from the other constraint rows and the indicator row.
Column $k$ becomes a unit column, while the unchanged basic columns stay unit
columns. The nonzero pivot also guarantees that the replacement basis is
nonsingular. Row numbers are \emph{not} variable indices: row $r$ is labelled
$x_{j_r}$ before the pivot and $x_k$ afterwards.

\section{The algorithm and degenerate steps}\label{sec:algorithm}
\begin{enumerate}[leftmargin=*]
  \item Start from a feasible basis and form $\bar A$, $\bar b$, $q$, and $z_0$.
  \item If all nonbasic indicators are nonnegative, return the current optimum.
  \item Select a nonbasic column $k$ with $q_k<0$.
  \item If $\bar a_k\leq0$, return an improving feasible ray and unboundedness.
  \item Apply the ratio test to positive entries of $\bar a_k$; choose a leaving row.
  \item Pivot, update the basis and objective row, and return to step 2.
\end{enumerate}
Without degeneracy each pivot strictly increases the objective, so no basis
can repeat; the finite number of bases implies termination. With degeneracy,
$\theta^*=0$ can change the basis without moving the point or increasing the
objective. Some pivot rules may cycle.

\textbf{Bland's rule} prevents cycling in exact arithmetic: choose the eligible
entering variable with the smallest \emph{original variable index}; among
minimum-ratio ties choose the leaving basic variable with the smallest original
index. Choosing only the topmost row in a tie is not this rule. Dantzig's rule
with arbitrary tie breaking has no general finite-termination guarantee.

\begin{figure}[H]
  \centering
  \begin{tikzpicture}[x=3cm,y=3cm,every node/.style={font=\small}]
    \filldraw[model region] (0,0)--(1,0)--(0,1)--cycle;
    \draw[model axis] (-0.08,0)--(1.45,0) node[right] {$x_1$};
    \draw[model axis] (0,-0.08)--(0,1.25) node[above] {$x_2$};
    \draw[modelorange,dashed] (1,-0.13)--(1,0.85);
    \node[model point] (v) at (1,0) {};
    \node[above right,align=left] at (1.04,0.16)
    {Same point $(1,0)$\\$B=(1,4,5)$\\$B'=(1,3,5)$};
    \node[above,rotate=-45] at (0.53,0.47) {$x_1+x_2=1$};
    \node[right] at (1,0.8) {$x_1=1$};
    \node[below,align=center] at (0.5,-0.23)
    {$x_3=1-x_1-x_2,\quad x_4=1-x_1,\quad x_5=1-x_2$};
  \end{tikzpicture}
  \caption{Degeneracy for $x_1+x_2\leq1$, $x_1\leq1$, $x_2\leq1$, $x\geq0$.
    At $(1,0)$ three boundary conditions meet. Both indicated bases give
    $x=(1,0,0,0,1)$; exchanging the zero slack changes the basis only.}
  \label{fig:degeneracy}
\end{figure}
Floating-point implementations also require scale-aware feasibility and
indicator tolerances. A theoretical zero may appear as a tiny signed number;
exact-arithmetic pivot logic cannot simply be applied without numerical care.

\clearpage
\section{Worked example: production planning}\label{sec:production}
\begin{quote}
  \foreignlanguage{polish}{Rozwiązać podstawową metodą sympleks.}
\end{quote}
Solve by the basic simplex method:
\[
  \max(3x_1+2x_2),\qquad
  2x_1+x_2\leq10,\quad 3x_1+3x_2\leq24,\quad
  \tfrac32x_1\leq6,\quad x_1,x_2\geq0.
\]
Production is continuous; selling prices are 3 and 2 PLN per item, so the
objective is revenue. The three resource limits are in kg. Add slacks
$x_3,x_4,x_5$ in that order and use
\[
  A=\begin{bmatrix}2&1&1&0&0\\3&3&0&1&0\\3/2&0&0&0&1\end{bmatrix},
  \quad b=\begin{bmatrix}10\\24\\6\end{bmatrix},\quad
  c=\begin{bmatrix}3\\2\\0\\0\\0\end{bmatrix}.
\]
Throughout the following tableaus the bottom row is $q=z-c$, with the current
objective value in the rightmost cell. Constraint rows are labelled by their
basic variables; all unlisted variables are zero at the current BFS.

\subsection{Start: the slack basis}
The ordered basis is $B=(3,4,5)$.
\begin{center}
  \begin{tabular}{crrrrrr}\toprule
    Basic & $x_1$ & $x_2$ & $x_3$ & $x_4$ & $x_5$ & RHS\\\midrule
    $x_3$ & 2 & 1 & 1 & 0 & 0 & 10\\
    $x_4$ & 3 & 3 & 0 & 1 & 0 & 24\\
    $x_5$ & $\boxed{3/2}$ & 0 & 0 & 0 & 1 & 6\\\midrule
    $q$ & $-3$ & $-2$ & 0 & 0 & 0 & $z_0=0$\\\bottomrule
  \end{tabular}
\end{center}
Dantzig's rule selects $x_1$, because $-3<-2$. Ratios are
$10/2=5$, $24/3=8$, and $6/(3/2)=4$. Hence $x_5$ leaves,
$\theta^*=4$, and revenue increases by $3\cdot4=12$.

\subsection{First pivot}
The ordered basis is $B=(3,4,1)$.
Divide the third row by $3/2$, then eliminate column $x_1$ from the first two
rows and the indicator row. For example, $R_1\leftarrow R_1-2R_3$ gives
$x_2+x_3-\frac43x_5=2$.
\begin{center}
  \begin{tabular}{crrrrrr}\toprule
    Basic & $x_1$ & $x_2$ & $x_3$ & $x_4$ & $x_5$ & RHS\\\midrule
    $x_3$ & 0 & $\boxed{1}$ & 1 & 0 & $-4/3$ & 2\\
    $x_4$ & 0 & 3 & 0 & 1 & $-2$ & 12\\
    $x_1$ & 1 & 0 & 0 & 0 & $2/3$ & 4\\\midrule
    $q$ & 0 & $-2$ & 0 & 0 & 2 & $z_0=12$\\\bottomrule
  \end{tabular}
\end{center}
The production point is $(4,0)$. Enter $x_2$: its positive-column ratios
are $2/1=2$ and $12/3=4$, while the third row has coefficient zero and is
excluded. Thus $x_3$ leaves and revenue gains $2\cdot2=4$.

\Needspace{13\baselineskip}
\subsection{Second pivot}
The ordered basis is $B=(2,4,1)$.
\begin{center}
  \begin{tabular}{crrrrrr}\toprule
    Basic & $x_1$ & $x_2$ & $x_3$ & $x_4$ & $x_5$ & RHS\\\midrule
    $x_2$ & 0 & 1 & 1 & 0 & $-4/3$ & 2\\
    $x_4$ & 0 & 0 & $-3$ & 1 & $\boxed{2}$ & 6\\
    $x_1$ & 1 & 0 & 0 & 0 & $2/3$ & 4\\\midrule
    $q$ & 0 & 0 & 2 & 0 & $-2/3$ & $z_0=16$\\\bottomrule
  \end{tabular}
\end{center}
At $(4,2)$, enter $x_5$. A slack can enter the basis just like any other
nonbasic variable. Its column is $(-4/3,2,2/3)^T$, so only rows 2 and 3
limit the step. Their ratios are $6/2=3$ and $4/(2/3)=6$. The first row's
negative coefficient means $x_2$ increases, not that its ratio should be tested.
Thus $x_4$ leaves with step 3 and revenue gain $(2/3)\cdot3=2$.

\subsection{Third pivot and certificate}
The ordered basis is $B=(2,5,1)$.
\begin{center}
  \begin{tabular}{crrrrrr}\toprule
    Basic & $x_1$ & $x_2$ & $x_3$ & $x_4$ & $x_5$ & RHS\\\midrule
    $x_2$ & 0 & 1 & $-1$ & $2/3$ & 0 & 6\\
    $x_5$ & 0 & 0 & $-3/2$ & $1/2$ & 1 & 3\\
    $x_1$ & 1 & 0 & 1 & $-1/3$ & 0 & 2\\\midrule
    $q$ & 0 & 0 & 1 & $1/3$ & 0 & $z_0=18$\\\bottomrule
  \end{tabular}
\end{center}
All indicators are nonnegative. Therefore
\[
  x^*=(2,6,0,0,3)^T,\qquad z^*=18\text{ PLN}.
\]
The first two resources are exhausted, while 3 kg of resource 3 remain.
The objective identity is an explicit certificate:
\begin{equation}\label{eq:production-certificate}
  3x_1+2x_2=18-x_3-\tfrac13x_4\leq18.
\end{equation}
Both nonbasic indicators are strictly positive, so this optimizer is unique.
The same upper bound follows by adding the first resource constraint to
one third of the second.

\begin{figure}[H]
  \centering
  \begin{minipage}[c]{0.52\linewidth}
    \centering
    \begin{tikzpicture}[x=0.75cm,y=0.48cm,every node/.style={font=\small}]
      \filldraw[model region] (0,0)--(4,0)--(4,2)--(2,6)--(0,8)--cycle;
      \draw[model axis] (-0.2,0)--(5.5,0) node[right] {$x_1$};
      \draw[model axis] (0,-0.3)--(0,9) node[above] {$x_2$};
      \draw[model flow,very thick] (0,0)--(4,0)--(4,2)--(2,6);
      \draw[model objective] (0.7,7.95)--(4.5,2.25);
      \foreach \p in {(0,0),(4,0),(4,2),(2,6)} \node[model point] at \p {};
      \node[below left] at (0,0) {$z=0$};
      \node[below] at (4,0) {$z=12$};
      \node[right] at (4,2) {$z=16$};
      \node[above right] at (2,6) {$z=18$};
      \node[left] at (0,8) {$8$};
      \node[right,align=left] at (5,7)
      {$x_1$ enters\\then $x_2$\\then slack $x_5$};
    \end{tikzpicture}
  \end{minipage}\hfill
  \begin{minipage}[c]{0.46\linewidth}
    \centering
    \begin{tabular}{ccrr}\toprule
      In & Out & Step & Revenue\\\midrule
      $x_1$ & $x_5$ & 4 & 12\\
      $x_2$ & $x_3$ & 2 & 16\\
      $x_5$ & $x_4$ & 3 & 18\\\bottomrule
    \end{tabular}
    \par\smallskip
    {\small Revenue in PLN.}
  \end{minipage}
  \caption{The tableau sequence follows three edges of the production polygon.
    The final dashed line is $3x_1+2x_2=18$. Each nondegenerate pivot moves
    to a distinct adjacent vertex. The table summarizes the entering and
    leaving variables, steps, and resulting revenue.}
  \label{fig:production-path}
\end{figure}

\section{Worked example: an unbounded objective}\label{sec:unbounded}
Consider
\[
  \max(3x_1+x_2-3x_3-x_4),\quad
  \begin{cases}
    x_1+2x_2-x_3-x_4\leq4,    \\
    2x_1-2x_2+3x_3-3x_4\leq9, \\
    x_1,x_2,x_3,x_4\geq0.
  \end{cases}
\]
Add slacks $x_5,x_6$. The initial basis is $(5,6)$ and the tableau is
\begin{center}
  \begin{tabular}{crrrrrrr}\toprule
    Basic & $x_1$ & $x_2$ & $x_3$ & $x_4$ & $x_5$ & $x_6$ & RHS\\\midrule
    $x_5$ & $\boxed{1}$ & 2 & $-1$ & $-1$ & 1 & 0 & 4\\
    $x_6$ & 2 & $-2$ & 3 & $-3$ & 0 & 1 & 9\\\midrule
    $q$ & $-3$ & $-1$ & 3 & 1 & 0 & 0 & $z_0=0$\\\bottomrule
  \end{tabular}
\end{center}
Enter $x_1$. The ratios $4$ and $9/2$ select the first row, so $x_5$ leaves.
After subtracting twice the first row from the second:
\begin{center}
  \begin{tabular}{crrrrrrr}\toprule
    Basic & $x_1$ & $x_2$ & $x_3$ & $x_4$ & $x_5$ & $x_6$ & RHS\\\midrule
    $x_1$ & 1 & 2 & $-1$ & $-1$ & 1 & 0 & 4\\
    $x_6$ & 0 & $-6$ & 5 & $-1$ & $-2$ & 1 & 1\\\midrule
    $q$ & 0 & 5 & 0 & $-2$ & 3 & 0 & $z_0=12$\\\bottomrule
  \end{tabular}
\end{center}
Now $q_4=-2$ and the $x_4$ column is $(-1,-1)^T$. There is no positive
entry and hence no leaving row. Set $x_4=\theta$, keeping $x_2=x_3=x_5=0$.
The equations give
\begin{equation}\label{eq:unbounded-ray}
  x(\theta)=(4+\theta,0,0,\theta,0,1+\theta)^T,
  \quad d=(1,0,0,1,0,1)^T,\quad \theta\geq0.
\end{equation}
Check directly: $Ad=0$, $d\geq0$, and $c^Td=2$. The original constraint
left-hand sides are $4$ and $8-\theta\leq9$, while the objective is
$12+2\theta\to+\infty$. Thus unboundedness comes with a verifiable feasible
ray, rather than just a failed ratio test.

This direction also spans an \emph{extreme ray} of the standard-form recession
cone $\{v:Av=0,v\geq0\}$. To see this, suppose $d=r+s$ with $r,s$ in that
cone. Since $d_2=d_3=d_5=0$, nonnegativity forces the same coordinates of
$r,s$ to be zero. The homogeneous equations then give $r_1=r_4$ and
$r_6=r_4$, so $r=r_4d$; the same holds for $s$. The decomposition can only
use parallel nonnegative multiples of $d$.

\begin{figure}[H]
  \centering
  \begin{tikzpicture}[x=0.8cm,y=0.8cm,every node/.style={font=\small}]
    \fill[modelblue!10] (0,0)--(4,0)--(7,3)--(7,4)--(0,4)--cycle;
    \draw[model axis] (0,0)--(7.8,0) node[right] {$x_1$};
    \draw[model axis] (0,0)--(0,4.5) node[above] {$x_4$};
    \draw[modelblue] (4,0)--(7.3,3.3);
    \draw[model flow,very thick] (4,0)--(6.7,2.7);
    \node[model point] at (4,0) {};
    \node[below] at (4,0) {$(4,0)$};
    \node[above left,align=center] at (5.5,1.5)
    {$(4+\theta,\theta)$\\$z=12+2\theta$};
    \node[left] at (0,3) {$3$};
    \node[right] at (7.1,3.8) {feasible slice};
  \end{tikzpicture}
  \caption{Slice $x_2=x_3=0$ in the original four-variable problem. The shaded
    feasible region continues upward and to the right; the arrow lies on
    $x_1-x_4=4$ and gives increasing objective.}
  \label{fig:unbounded-ray}
\end{figure}

\Needspace{18\baselineskip}
\section{Finding an initial feasible basis}\label{sec:initialization}
A simplex step preserves feasibility; it cannot repair an infeasible starting
basis. Before using the primal algorithm, convert the constraints to equalities
and establish a BFS.
\begin{center}
  \begin{tabular}{>{\raggedright\arraybackslash}p{0.23\linewidth}>{\raggedright\arraybackslash}p{0.32\linewidth}>{\raggedright\arraybackslash}p{0.35\linewidth}}\toprule
    Constraint or situation & Equality form & Initial basis implication\\\midrule
    $a^Tx\leq b$, $b\geq0$ & $a^Tx+s=b$, $s\geq0$ & Slack column $+1$ can start basic with $s=b$.\\
    $a^Tx\geq b$, $b\geq0$ & $a^Tx-s=b$, $s\geq0$ & Surplus column $-1$ alone gives $s=-b$; generally add an artificial column.\\
    $a^Tx=b$ & Equality already present & Reuse a feasible basis if available; otherwise add an artificial column.\\
    Equality with $b<0$ & Multiply the entire row by $-1$ & Normalize RHS before introducing a positive artificial column.\\
    Inequality with $b<0$ & Multiply by $-1$ and reverse inequality & Then choose slack or surplus for the new direction.\\\bottomrule
  \end{tabular}
\end{center}
Slacks and surplus variables describe the original constraint. Artificial
variables are temporary additions used to find feasibility; they must be zero
in a solution of the original problem.

For $Ax\leq b$, $b\geq0$, slacks give the immediate basis $x=0$, $s=b$.
More generally, identity columns already present in an equality system can form
an initial basis when their RHS values are nonnegative. An identity block with
a negative basic value is not a feasible basis. Zero RHS values are allowed but
produce degeneracy.

\subsection{Phase I: a feasibility problem}\label{sec:phase-one}
After normalizing equality rows to $b\geq0$, a universally available augmented
system is
\[
  Ax+a=b,\qquad x\geq0,\quad a\geq0,
\]
where each row has an artificial variable. It has the BFS $x=0$, $a=b$.
If some usable basic columns already exist, add artificials only to the rows
that need them. Phase I solves
\begin{equation}\label{eq:phase-one}
  \min w=\sum_i a_i\quad\text{or, equivalently,}\quad\max(-w).
\end{equation}
Always build its indicator row from its actual basic objective coefficients;
the artificial variables may already be basic and cannot be treated as if their
objective contribution were absent.

\begin{proposition}[Phase I test]
  The original equality/nonnegativity system is feasible if and only if the
  Phase I optimum is $w^*=0$.
\end{proposition}
\begin{proof}
  A feasible original $x$ gives an augmented feasible point $(x,0)$ with $w=0$.
  Since $w\geq0$, this is the minimum. Conversely, $w^*=0$ and $a_i\geq0$
  force every artificial variable to vanish, so $Ax=b$, $x\geq0$.
\end{proof}
Phase I is feasible by construction and its objective is bounded below by zero,
so it has a finite attained minimum. A positive \emph{intermediate} $w$ says
nothing about infeasibility; only a positive \emph{optimal} $w^*$ does.

\begin{figure}[H]
  \centering
  \begin{tikzpicture}[every node/.style={font=\small},
      box/.style={draw=modelblue,rounded corners,align=center,inner sep=6pt}]
    \node[box] (start) at (0,0) {Normalize RHS; add needed artificials\\solve Phase I};
    \node[box] (no) at (-3.5,-1.7) {$w^*>0$\\original problem infeasible};
    \node[box] (yes) at (3,-1.7) {$w^*=0$\\original problem feasible};
    \node[box] (clean) at (3,-3.2) {Remove zero basic artificials\\and any redundant rows};
    \node[box] (two) at (3,-4.6) {Delete artificial columns\\recompute original objective; Phase II};
    \draw[model flow] (start)--(no);
    \draw[model flow] (start)--(yes);
    \draw[model flow] (yes)--(clean);
    \draw[model flow] (clean)--(two);
  \end{tikzpicture}
  \caption{Two-phase initialization. Phase I establishes feasibility;
    Phase II optimizes the original objective.}
  \label{fig:phase-flow}
\end{figure}

\subsection{Removing artificial variables from a zero optimum}
At $w^*=0$, a zero artificial variable may remain \emph{basic}. Removing its
column immediately would leave an incomplete basis. Inspect its transformed
row restricted to original nonbasic columns:
\begin{itemize}[leftmargin=*]
  \item If a coefficient is nonzero, pivot that original variable into the row.
        The RHS is zero, so the point stays unchanged. For this cleanup exchange
        the coefficient may be positive or negative: this is a basis repair at
        zero, not an objective-improving positive-step ratio test.
  \item If every original coefficient in that row is zero, its original
        equality is $0=0$ after transformation. Delete that redundant row and its
        artificial basic variable.
\end{itemize}
Under the full-row-rank assumption on the original $A$, such a redundant row
cannot occur; it matters when dependent rows have been retained. For instance,
$x_1=1$ and $2x_1=2$ reduce to one independent equality. In contrast,
$x_1=1$ and $2x_1=3$ are inconsistent, not a redundant pair.
After cleanup, delete artificial columns and recompute $q$ and $z_0$ using the
\emph{original} objective and the retained basis. Phase I indicators cannot
certify optimality for Phase II.

\subsection{A small infeasible example}
For $x\geq2$ and $x\leq1$, introduce surplus $s$, slack $t$, and artificials:
\[
  x-s+a_1=2,\qquad x+t+a_2=1,\qquad x,s,t,a_1,a_2\geq0.
\]
Subtracting the second equation from the first gives
$a_1=1+s+t+a_2$. Therefore $w=a_1+a_2\geq1$. The augmented point
$(x,s,t,a_1,a_2)=(1,0,0,1,0)$ attains $w=1$, so $w^*=1>0$ and the
original problem is infeasible. This proves the Phase I result without relying
on a numerical solver status.

\clearpage
\section{Worked example: two different Phase I starts}\label{sec:two-phase}
Solve
\begin{equation}\label{eq:two-phase-original}
  \max(-x_1-2x_2-3x_3-4x_4),\qquad
  \begin{cases}
    x_1+x_2+x_3+x_4=1,     \\
    x_1+x_3-3x_4=\tfrac12, \\
    x_1,x_2,x_3,x_4\geq0.
  \end{cases}
\end{equation}
The coefficient matrix has rank 2, but its original columns do not provide an
immediate feasible identity basis. Compare an all-artificial start with a start
that reuses the existing column $a_2=(1,0)^T$.

\subsection{Start A: two artificial variables}
Introduce artificial variables $x_5,x_6$ and maximize $-x_5-x_6$ in Phase I:
\[
  x_1+x_2+x_3+x_4+x_5=1,\qquad
  x_1+x_3-3x_4+x_6=\tfrac12.
\]
The basis $(5,6)$ gives $x_5=1$, $x_6=1/2$, and $w=3/2$.
Use the Phase I coefficients $c^{\mathrm I}=(0,0,0,0,-1,-1)^T$:
\begin{center}
  \begin{tabular}{crrrrrrr}\toprule
    Basic & $x_1$ & $x_2$ & $x_3$ & $x_4$ & $x_5$ & $x_6$ & RHS\\\midrule
    $x_5$ & 1 & 1 & 1 & 1 & 1 & 0 & 1\\
    $x_6$ & $\boxed{1}$ & 0 & 1 & $-3$ & 0 & 1 & $1/2$\\\midrule
    $q^{\mathrm I}$ & $-2$ & $-1$ & $-2$ & 2 & 0 & 0 & $z_{\mathrm I}=-3/2$\\\bottomrule
  \end{tabular}
\end{center}
The most negative indicators tie for $x_1$ and $x_3$; select $x_1$ by smaller
index. Ratios $1$ and $1/2$ select row 2, so $x_6$ leaves.
\begin{center}
  \begin{tabular}{crrrrrrr}\toprule
    Basic & $x_1$ & $x_2$ & $x_3$ & $x_4$ & $x_5$ & $x_6$ & RHS\\\midrule
    $x_5$ & 0 & 1 & 0 & $\boxed{4}$ & 1 & $-1$ & $1/2$\\
    $x_1$ & 1 & 0 & 1 & $-3$ & 0 & 1 & $1/2$\\\midrule
    $q^{\mathrm I}$ & 0 & $-1$ & 0 & $-4$ & 0 & 2 & $z_{\mathrm I}=-1/2$\\\bottomrule
  \end{tabular}
\end{center}
Now $x_4$ enters and only row 1 has a positive coefficient. The step is
$(1/2)/4=1/8$, and $x_5$ leaves. The Phase I objective increases by
$4\cdot(1/8)=1/2$ to zero.
\begin{center}
  \begin{tabular}{crrrrrrr}\toprule
    Basic & $x_1$ & $x_2$ & $x_3$ & $x_4$ & $x_5$ & $x_6$ & RHS\\\midrule
    $x_4$ & 0 & $1/4$ & 0 & 1 & $1/4$ & $-1/4$ & $1/8$\\
    $x_1$ & 1 & $3/4$ & 1 & 0 & $3/4$ & $1/4$ & $7/8$\\\midrule
    $q^{\mathrm I}$ & 0 & 0 & 0 & 0 & 1 & 1 & $z_{\mathrm I}=0$\\\bottomrule
  \end{tabular}
\end{center}
Both artificials are nonbasic and zero. Delete columns 5 and 6; retain
$B=(4,1)$. Recompute the objective using $c^{\mathrm{II}}=(-1,-2,-3,-4)^T$:
\begin{center}
  \begin{tabular}{crrrrr}\toprule
    Basic & $x_1$ & $x_2$ & $x_3$ & $x_4$ & RHS\\\midrule
    $x_4$ & 0 & $1/4$ & 0 & 1 & $1/8$\\
    $x_1$ & 1 & $3/4$ & 1 & 0 & $7/8$\\\midrule
    $q^{\mathrm{II}}$ & 0 & $1/4$ & 2 & 0 & $z_{\mathrm{II}}=-11/8$\\\bottomrule
  \end{tabular}
\end{center}
All Phase II indicators are nonnegative, so this start has already reached the
original optimum. Phase I happened to find a basis optimal for Phase II;
this is not guaranteed in general.

\subsection{Start B: reuse an original basic column}
Keep $x_2$ basic in row 1 and introduce only $x_5=a$ in row 2:
\[
  x_1+x_2+x_3+x_4=1,\qquad x_1+x_3-3x_4+x_5=\tfrac12.
\]
Start with $B=(2,5)$ and maximize $-x_5$.
\begin{center}
  \begin{tabular}{crrrrrr}\toprule
    Basic & $x_1$ & $x_2$ & $x_3$ & $x_4$ & $x_5$ & RHS\\\midrule
    $x_2$ & 1 & 1 & 1 & 1 & 0 & 1\\
    $x_5$ & $\boxed{1}$ & 0 & 1 & $-3$ & 1 & $1/2$\\\midrule
    $q^{\mathrm I}$ & $-1$ & 0 & $-1$ & 3 & 0 & $z_{\mathrm I}=-1/2$\\\bottomrule
  \end{tabular}
\end{center}
Enter $x_1$ and remove $x_5$ after step $1/2$:
\begin{center}
  \begin{tabular}{crrrrrr}\toprule
    Basic & $x_1$ & $x_2$ & $x_3$ & $x_4$ & $x_5$ & RHS\\\midrule
    $x_2$ & 0 & 1 & 0 & 4 & $-1$ & $1/2$\\
    $x_1$ & 1 & 0 & 1 & $-3$ & 1 & $1/2$\\\midrule
    $q^{\mathrm I}$ & 0 & 0 & 0 & 0 & 1 & $z_{\mathrm I}=0$\\\bottomrule
  \end{tabular}
\end{center}
Delete column 5. The feasible original point is $(1/2,1/2,0,0)^T$, but
recomputing its \emph{original} objective gives
\begin{center}
  \begin{tabular}{crrrrr}\toprule
    Basic & $x_1$ & $x_2$ & $x_3$ & $x_4$ & RHS\\\midrule
    $x_2$ & 0 & 1 & 0 & $\boxed{4}$ & $1/2$\\
    $x_1$ & 1 & 0 & 1 & $-3$ & $1/2$\\\midrule
    $q^{\mathrm{II}}$ & 0 & 0 & 2 & $-1$ & $z_{\mathrm{II}}=-3/2$\\\bottomrule
  \end{tabular}
\end{center}
The negative indicator for $x_4$ requires a Phase II pivot. Enter $x_4$,
remove $x_2$, and take step $1/8$. Dividing row 1 by 4 and adding three
copies to row 2 produces exactly the final Phase II tableau from Start A.
Therefore both starts yield
\begin{equation}\label{eq:two-phase-result}
  x^*=(7/8,0,0,1/8)^T,\qquad z^*=-11/8.
\end{equation}
Its indicator identity gives a full hand certificate:
\begin{equation}\label{eq:two-phase-certificate}
  -x_1-2x_2-3x_3-4x_4=-\tfrac{11}{8}-\tfrac14x_2-2x_3
  \leq-\tfrac{11}{8}.
\end{equation}
The two nonbasic indicators are positive, so the optimum is unique.

\subsection{Geometry of the feasible equality system}
Eliminate $x_1,x_2$ from Equation~\eqref{eq:two-phase-original}:
\[
  x_1=\tfrac12-x_3+3x_4,\qquad x_2=\tfrac12-4x_4.
\]
Nonnegativity gives
$0\leq x_4\leq1/8$ and $0\leq x_3\leq1/2+3x_4$.
These two coordinates parameterize the whole feasible set, and
\[
  z=-\tfrac32+x_4-2x_3.
\]
The maximum is at the largest allowed $x_4$ and smallest allowed $x_3$.
This independent calculation confirms the tableau result.
\begin{figure}[H]
  \centering
  \begin{tikzpicture}[x=24cm,y=4cm,every node/.style={font=\small}]
    \filldraw[model region] (0,0)--(0.125,0)--(0.125,0.875)--(0,0.5)--cycle;
    \draw[model axis] (-0.005,0)--(0.18,0) node[right] {$x_4$};
    \draw[model axis] (0,-0.04)--(0,1.04) node[above] {$x_3$};
    \draw[model flow,very thick] (0,0)--(0.125,0);
    \node[model point] at (0,0) {};
    \node[model point] at (0.125,0) {};
    \node[below left,align=right] at (0,0)
    {Start B after Phase I\\$z=-3/2$};
    \node[below right,align=left] at (0.125,0)
    {Both starts finish here\\$z=-11/8$};
    \node[left] at (0,0.5) {$1/2$};
    \node[right] at (0.125,0.875) {$(1/8,7/8)$};
    \node[above,rotate=27] at (0.065,0.695) {$x_3=1/2+3x_4$};
    \node[right,align=left] at (0.18,0.55)
    {Start A reaches the optimum\\during Phase I.\\Start B needs the shown\\Phase II edge move.};
  \end{tikzpicture}
  \caption{The four-variable equality system is a quadrilateral in coordinates
    $(x_4,x_3)$. Its lower edge has increasing objective toward the unique optimum.}
  \label{fig:two-phase-geometry}
\end{figure}
\begin{center}
  \begin{tabular}{lrrl}\toprule
    Initialization & Artificial variables & Phase I pivots & Phase II pivots\\\midrule
    Start A & 2 & 2 & 0\\
    Start B & 1 & 1 & 1\\\bottomrule
  \end{tabular}
\end{center}
Reusing an existing basis column saves an artificial variable here. It does not
imply that every such reuse minimizes the total number of simplex pivots.

\section{The Big M method and its limitations}\label{sec:big-m}
After normalizing the RHS and introducing needed artificials, the maximization
penalty formulation is
\begin{equation}\label{eq:big-m}
  \max\bigl(c^Tx-M\textstyle\sum_i a_i\bigr),\qquad
  Ax+a=b,\quad x,a\geq0.
\end{equation}
Here the displayed form adds one artificial per row; a mixed starting basis
can use fewer. Positive artificial values are penalized, but $M$ must dominate
the competing original objective for the intended feasibility priority.
A rigorous symbolic interpretation compares objective expressions
lexicographically: first minimize $\sum a_i$, then optimize $c^Tx$ among
points with the minimum artificial sum. This recovers the two-phase priority.
When comparing $\alpha M+\beta$ symbolically, compare $\alpha$ first and then
$\beta$; an arbitrary large floating-point constant is not the same operation.

\subsection{A symbolic Big M calculation for the equality example}
Apply the penalty to Equation~\eqref{eq:two-phase-original}, using the same two
artificial columns as Start A. The combined coefficient vector is
$c^M=(-1,-2,-3,-4,-M,-M)^T$. Its initial basis $(5,6)$ gives
\[
  q^M=(1-2M,\ 2-M,\ 3-2M,\ 4+2M,\ 0,\ 0)^T,
  \qquad z_M=-3M/2.
\]
For symbolic dominant $M$, $x_1$ enters and $x_6$ leaves after step $1/2$.
The constraint rows are exactly the first transformed rows of Start A;
only their objective coefficients differ. Recompute with the combined
objective at $B=(5,1)$:
\[
  q^M=(0,\ 2-M,\ 2,\ 7-4M,\ 0,\ 2M-1)^T,
  \qquad z_M=-(M+1)/2.
\]
Now $x_4$ enters and $x_5$ leaves after step $1/8$. At $B=(4,1)$,
\[
  q^M=(0,\ 1/4,\ 2,\ 0,\ M-7/4,\ M+3/4)^T,
  \qquad z_M=-11/8.
\]
Both artificials are zero and every indicator is nonnegative for dominant $M$.
This certifies the same original optimum $(7/8,0,0,1/8)^T$ as the two-phase
method. For this \emph{specific} problem, the same Dantzig pivot sequence and
strict nonbasic indicator certificate also work with any finite $M>7/4$.
For example $M=10$ gives penalty-objective values $-15$, $-11/2$, and $-11/8$.
The threshold depends on the data; it is not a rule for other LPs.

\subsection{A hand counterexample to an insufficient penalty}
The original problem $\max(-Kx)$ subject to $x=1$, $x\geq0$ is feasible,
with unique value $-K$. Its augmented problem is
\[
  \max(-Kx-Ma),\qquad x+a=1,\quad x,a\geq0,
  \qquad K>M>0.
\]
Substitute $a=1-x$:
\[
  -Kx-M(1-x)=-M+(M-K)x,\qquad 0\leq x\leq1.
\]
The maximum is at $x=0,a=1$, with value $-M>-K$. A positive artificial
variable at this penalized optimum does \emph{not} prove the original system
infeasible. For $K=100$, $M=10$, the penalized optimum is $-10$ while the
original feasible point has value $-100$.
\begin{figure}[H]
  \centering
  \begin{tikzpicture}[x=4cm,y=0.035cm,every node/.style={font=\small}]
    \draw[model axis] (0,0)--(1.4,0) node[right] {$x$};
    \draw[model axis] (0,-110)--(0,8) node[above] {penalized objective};
    \draw[modelorange,very thick] (0,-10)--(1,-100);
    \node[model point] at (0,-10) {};
    \node[model point] at (1,-100) {};
    \node[right,align=left] at (0.16,-10) {$a=1$, value $-10$\\penalized optimum};
    \node[right,align=left] at (1,-100) {$a=0$, value $-100$\\original feasible point};
    \draw[dashed,black!40] (1,0)--(1,-100);
    \node[above] at (1,0) {$1$};
  \end{tikzpicture}
  \caption{With $K=100$ and $M=10$, the penalty rewards keeping an artificial
    variable because avoiding the original cost is more valuable.}
  \label{fig:big-m}
\end{figure}

\subsection{What a penalized solve can certify}
A \emph{finite optimal} augmented solution with $a=0$ certifies an original
optimum even for finite $M$: every original feasible point also belongs to the
augmented system and has the same objective there. In contrast, interpreting
positive artificials requires a justified feasibility priority.

An unbounded penalized problem with a positive artificial variable on a ray
also does not alone prove original infeasibility. For example,
$x_1-x_2+a=1$ and objective $x_1-Ma$ admit the augmented ray
$(x_1,x_2,a)=(\theta,\theta,1)$ with increasing objective. The original
$a=0$ system is nevertheless feasible and unbounded, using
$(x_1,x_2)=(1+\theta,\theta)$.
To certify original unboundedness, establish an original feasible point and
an improving ray whose artificial components are zero.

\begin{center}
  \begin{tabular}{>{\raggedright\arraybackslash}p{0.22\linewidth}>{\raggedright\arraybackslash}p{0.33\linewidth}>{\raggedright\arraybackslash}p{0.35\linewidth}}\toprule
    Aspect & Two-phase method & Big M method\\\midrule
    Feasibility priority & Separate minimization of artificial sum & Dominant symbolic penalty or justified numerical $M$\\
    Objective rows & Recompute the original row after Phase I & Original and penalty coefficients share a row\\
    Infeasibility test & $w^*>0$ at Phase I optimum & Positive artificials require a valid dominance argument\\
    Numerical issue & Ordinary feasibility/optimality tolerances & Large coefficients can also amplify scaling and cancellation problems\\\bottomrule
  \end{tabular}
\end{center}
For hand calculations, both formulations explain artificial variables. For a
transparent numerical demonstration, two separate phases avoid inventing a
penalty size. A simplex tableau is a teaching representation; efficient solvers
usually maintain basis factorizations instead of repeatedly forming full inverses.

\section{From a full tableau to revised simplex}\label{sec:revised}
The algebraic inverse $A_B^{-1}$ explains a basis transformation. An
implementation can obtain the same information by solving linear systems:
\[
  A_Bx_B=b,\qquad A_B^Ty=c_B,\qquad q_j=a_j^Ty-c_j.
\]
After selecting an entering $k$ with $q_k<0$, solve $A_Bu=a_k$.
Apply the same ratio test to $(x_B)_i/u_i$ over rows with $u_i>0$.
If none exist, the direction certifies unboundedness. This is \emph{revised simplex}:
it computes the needed quantities without storing the entire transformed
matrix $A_B^{-1}A$. With identical exact pivot choices, it follows the same
bases as the tableau method and has the same degeneracy issues.

\begin{figure}[H]
  \centering
  \begin{tikzpicture}[font=\small,
      box/.style={draw=modelblue,rounded corners,align=center,text width=3.5cm,inner sep=6pt}]
    \node[box] (factor) at (0,0) {Factor the current $A_B$\\solve for $x_B$ and $y$};
    \node[box] (price) at (4.9,0) {Price nonbasic\\columns\\$q_j=a_j^Ty-c_j$};
    \node[box] (step) at (9.8,0) {Solve $A_Bu=a_k$\\ratio test\\exchange column};
    \draw[model flow] (factor)--(price);
    \draw[model flow] (price)--(step);
    \draw[model flow] (step.south)--(9.8,-1.3)--(0,-1.3)--(factor.south);
    \node[fill=white] at (4.9,-1.3) {Update or refactor the new basis};
  \end{tikzpicture}
  \caption{The revised formulation retains the mathematical basis exchange.
    Pricing selects an improving column or establishes that all indicators
    are nonnegative.}\label{fig:revised-simplex}
\end{figure}
Sparse factorizations and their updates make this useful for large models;
periodic refactorization limits accumulated numerical error. Presolve can
remove fixed variables and redundant rows, but solutions and certificates must
be mapped back to the original model. A small equation residual alone does not
guarantee accurate coordinates when $A_B$ is ill-conditioned.
Scale variables and rows sensibly, distinguish numerical tolerances from exact
sign tests, and verify original feasibility as well as dual indicators.

\section{Checking the hand calculations in MATLAB}\label{sec:matlab}
The companion script is
\href{../../matlab/exercises/lecture_03_checks.m}{\texttt{matlab/exercises/lecture\_03\_checks.m}}.
Run it from the repository root:
\begin{lstlisting}
run('matlab/exercises/lecture_03_checks.m')
\end{lstlisting}
It uses base MATLAB linear solves and assertions, without Optimization Toolbox.
It checks the stated bases and pivot choices rather than implementing a general
purpose LP solver. The same matrices and objective coefficients are used in the
hand calculations and the script.

\Needspace{10\baselineskip}
A basis calculation is just two linear solves:
\begin{lstlisting}
AB = A(:, B);
Abar = AB \ A;
bbar = AB \ b;
q = Abar' * c(B) - c;
x = zeros(size(A, 2), 1);
x(B) = bbar;
value = c' * x;
\end{lstlisting}
Here $B$ is an ordered row vector of indices; $b,c,x,q$ are column vectors.
The ratio test deliberately discards zero and negative column coefficients:
\begin{lstlisting}
u = Abar(:, entering);
rows = find(u > tol);
if isempty(rows)
    direction = zeros(size(A, 2), 1);
    direction(B) = -u;
    direction(entering) = 1;
else
    ratios = bbar(rows) ./ u(rows);
    step = min(ratios);
end
\end{lstlisting}
This fragment illustrates the algebra; the companion script additionally checks
entering indicators, leaving choices, feasibility, tableau row operations, and
agreement with the next independently recomputed basis. It uses an absolute
tolerance $10^{-10}$ appropriate for these small, modestly scaled inputs.
A generic solver needs scale-aware tolerances and a strategy for ill-conditioned
bases. Exact fractions in the notes are not rounded numerical results.

\begin{center}
  \begin{tabular}{>{\raggedright\arraybackslash}p{0.25\linewidth}>{\raggedright\arraybackslash}p{0.38\linewidth}>{\raggedright\arraybackslash}p{0.27\linewidth}}\toprule
    Check & Expected result & Interpretation\\\midrule
    Production basis sequence & Revenue $0,12,16,18$ PLN & Three improving pivots\\
    Production final point & $(2,6,0,0,3)^T$ & Slacks in kg\\
    Unbounded example & $Ad=0$, $d\geq0$, $c^Td=2$ & Objective $12+2\theta$\\
    Two-phase Start A & $w:3/2,1/2,0$ & Two feasibility pivots\\
    Two-phase Start B & $w:1/2,0$ & One feasibility pivot\\
    Original equality problem & $(7/8,0,0,1/8)^T$, $z=-11/8$ & Same optimum\\
    Infeasible inequalities & $w^*=1$ & Infeasible\\
    Big M equality solve & $M=10$, $z=-11/8$, $a=0$ & Original optimum\\
    Finite penalty example & $x=0,a=1$, value $-10$ & $M=10<K=100$\\\bottomrule
  \end{tabular}
\end{center}

\Needspace{20\baselineskip}
\section{Study questions and common mistakes}\label{sec:review}
Before doing another tableau, check these distinctions:
\begin{center}
  \small
  \begin{tabular}{>{\raggedright\arraybackslash}p{0.35\linewidth}>{\raggedright\arraybackslash}p{0.57\linewidth}}\toprule
    Mistake & Correct check\\\midrule
    Every basic solution is feasible & Verify both $Ax=b$ and $x\geq0$.\\
    A zero variable always means degeneracy & Only a zero \emph{basic} coordinate establishes degeneracy.\\
    An unbounded set means an unbounded LP & Exhibit a feasible ray with \emph{positive} objective slope.\\
    Test every ratio in the entering column & Include only strictly positive coefficients.\\
    A zero indicator proves another optimum & Verify a positive feasible step or otherwise establish a distinct optimal point.\\
    A simplex pivot always moves the point & A zero step can change only the basis.\\
    A positive Phase I value proves infeasibility & Complete Phase I to its optimum first.\\
    Delete all artificial columns immediately & First replace zero basic artificials or remove redundant rows.\\
    Reuse Phase I indicators in Phase II & Recompute with the original objective coefficients.\\
    Any large $M$ is enough & Justify dominance; positive artificial values alone do not certify infeasibility.\\\bottomrule
  \end{tabular}
\end{center}
\begin{enumerate}[leftmargin=*,itemsep=0pt,topsep=4pt]
  \item In the production example, compute the initial pivot if $x_2$ enters
        first. How does its vertex path differ from Dantzig's path?
  \item Starting at the final production basis, explain why increasing either
        nonbasic slack cannot improve revenue.
  \item In the unbounded example, verify the entire ray in the original
        inequalities, without using the transformed tableau.
  \item In Start B of the two-phase example, derive $q_4^{\mathrm{II}}=-1$
        directly from $c_B^T\bar a_4-c_4$.
  \item For the finite penalty counterexample, characterize all optimizers
        when $M<K$, $M=K$, and $M>K$.
\end{enumerate}
\paragraph{Answers for checking.}
(1) Ratios $10$ and $8$ give step 8, so $x_4$ leaves and the point is $(0,8)$;
then $x_1$ enters with step 2 and $x_3$ leaves, reaching $(2,6)$.
(2) Equation~\eqref{eq:production-certificate} subtracts $x_3+x_4/3$ from 18.
(3) The two left-hand sides are $4$ and $8-\theta$.
(4) $(-2)4+(-1)(-3)-(-4)=-1$.
(5) Respectively $(x,a)=(0,1)$, the whole segment $x+a=1$ with $x,a\geq0$,
and $(x,a)=(1,0)$.

\clearpage
\section{Sources and editorial remarks}\label{sec:sources}
\noindent\textbf{Primary source:} E. Pawelec, \texttt{s4\_w3.pdf},
PDF pages 1--12, printed slides 1--46.\\
\textbf{Source date:} Stored in the folder labelled 2020; PDF metadata records
creation on 2019-10-17, which does not establish the delivery date.\\
\textbf{Current course alignment:} \texttt{ProgMat\_w1.pdf}, pp.~32--39,
confirms simplex and initialization coverage. Pages~36--39 contain the same
equality example developed in Section~\ref{sec:two-phase}. Sheet~1, group~1
requires the methods practiced in \texttt{01-primal-simplex.tex}.\\
\textbf{Editorial scope:} The complete mathematical content is reorganized in
English. The worked tableaus retain the supplied data and both Phase I starts.
Proofs, counterexamples, study questions, redrawn figures, and the MATLAB
companion are additional explanations.

\begin{center}
  \begin{tabular}{lll}\toprule
    PDF pages & Printed slides & Material in these notes\\\midrule
    1 & 1--4 & Section~\ref{sec:problem}; geometric interpretation\\
    2--3 & 5--12 & Sections~\ref{sec:bases}--\ref{sec:canonical}; candidate solutions\\
    4--6 & 13--23 & Sections~\ref{sec:canonical}--\ref{sec:algorithm}; simplex derivation\\
    6--7 & 24--28 & Section~\ref{sec:production}; full production tableaus\\
    8 & 29--31 & Section~\ref{sec:unbounded}; unbounded example\\
    8--9 & 32--36 & Section~\ref{sec:initialization}; initialization and two phases\\
    10 & 37--38 & Section~\ref{sec:big-m}; penalty method\\
    10--11 & 39--44 & Section~\ref{sec:two-phase}; two initialization variants\\
    12 & 45--46 & Section~\ref{sec:row-operations}; pivot formulas\\\bottomrule
  \end{tabular}
\end{center}
The Polish prompt in Section~\ref{sec:production} is from slide 24.
Candidate vectors in Section~\ref{sec:bases} are from slides 8 and 12.
The unbounded ray is fully derived here from the example and question on
slides 29--31. All figures are editable explanatory TikZ redrawings; the ratio
test sketch, degeneracy example, and finite-penalty graph use added illustrative
data. The equations and tableaus on slides 39--44 were inspected visually to
confirm nonnegativity of all four original variables and both artificial-variable
conventions.

\paragraph{Current-source qualifications.}
In \texttt{ProgMat\_w1.pdf}, p.~34, an easily usable basic column should
have a positive entry in its pivot row and zero entries in the other rows;
the printed indices instead test other entries in a row. Normalize a positive
column before reading its basic value. The penalty claims on page~35
require the Big M qualifications in Section~\ref{sec:big-m}. The image-based
tableaus on pp.~37--39 were checked visually: the all-artificial path already
developed here has Phase I pivots entering $x_1$ and then $x_4$, and Phase II
needs no further pivot for that path. This does not imply that Phase II is
generally unnecessary.

\paragraph{Qualifications of the mathematical claims.}
Section~\ref{sec:canonical} distinguishes a sufficient indicator certificate
from necessity for a nondegenerate BFS, qualifying the equivalence stated on
slide 14. It also requires an actual positive feasible step to infer a distinct
optimizer from a zero indicator. Section~\ref{sec:initialization} explains
zero-basic-artificial cleanup and redundant rows explicitly.
Section~\ref{sec:big-m} qualifies the classification on slide 38:
an insufficient finite penalty and an augmented ray carrying positive artificial
values do not by themselves establish original infeasibility.

\paragraph{Supplementary checks.}
Section~\ref{sec:revised} and Figure~\ref{fig:revised-simplex} connect the
existing pivot equations to implementation. The terminology was checked against
\href{https://ocw.mit.edu/courses/6-251j-introduction-to-mathematical-programming-fall-2009/2ef2f1dd7045b5f29e5faea299fb0798_MIT6_251JF09_lec06.pdf}{MIT OCW,
  \emph{The Simplex Method II}, Lecture 6}, PDF pp.~2--3.
This is additional study material with unknown current-course alignment;
no new solver implementation or MATLAB execution is claimed for this section.

\href{https://ocw.mit.edu/courses/6-251j-introduction-to-mathematical-programming-fall-2009/4fdf1a06497288c24593f7eedef7576e_MIT6_251JF09_lec05.pdf}{MIT OCW 15.081J/6.251J, Lecture 5 (Fall 2009)},
PDF p.~2 and pp.~5--6, was consulted for the nondegeneracy qualification and
anti-cycling rule.
\href{https://ocw.mit.edu/courses/6-251j-introduction-to-mathematical-programming-fall-2009/12b1c7e885991696c294953b2497739c_MIT6_251JF09_lec07.pdf}{Lecture 7},
PDF pp.~2--4, was consulted for removal of zero basic artificials and the
transition to Phase II. These are supplementary references, not evidence of
current course coverage; the source inventory records external storage status.

\paragraph{Verification.}
All displayed worked-example tableaus were independently recomputed using
exact rational Gauss--Jordan elimination, with constraint residuals checked
exactly. Their objective and indicator rows agree with the hand derivations.
The installed R2026b static code analyzer checked the MATLAB companion with
no diagnostics. On 2026-10-04 the user ran the companion in a normal terminal
and supplied its output: all numerical assertions passed. The reported production
sequence, unbounded-ray slope, two-phase optimum, and infeasibility result agree
with Section~\ref{sec:matlab}. The run did not print a version; an earlier
user run reported MATLAB R2026b. The supplied terminal output is recorded in
\texttt{docs/matlab-verification.md}.
Polish terminology is maintained in \texttt{docs/glossary.md}.
\end{document}
